% ID: FIS-MAG-FAR-001
% Titolo: Significato fisico del segno meno nella legge di Lenz
% Disciplina: Fisica
% Area: Magnetismo
% Argomento: Faraday-Lenz
% Sottoargomento: Conservazione dell'energia e verso della corrente indotta
% Classe: Quinta liceo scientifico
% Difficolta: 3
% Tipo: quesito_teorico
% Risultato: soluzione da aggiungere
% Tempo_stimato: 7 min
% Competenze: interpretazione fisica, collegamento con la conservazione dell'energia, argomentazione
% Prerequisiti: flusso magnetico, induzione elettromagnetica, legge di Faraday
% Tag: fisica, magnetismo, faraday_lenz, legge_lenz, induzione
% Fonte: verifica_fisica_magnetismo_anonimizzata
% Autore: Riccardo Carli
% Licenza: CC-BY-SA-4.0

\begin{esercizio}
Spiega in breve il motivo fisico per cui nella legge di Faraday-Lenz compare
un segno meno davanti alla variazione del flusso del campo magnetico nel tempo.
\end{esercizio}

\begin{soluzione}
% Soluzione da aggiungere.
\end{soluzione}

